\documentclass[]{note}
\usepackage{CJKutf8}
\usepackage{amsmath}
\usepackage{amsfonts}
\usepackage{braket}
\usepackage[ruled,vlined]{algorithm2e}
\def\redd#1{\textcolor{red}{#1}}
\def\dom{\text{dom}}
\def\ri{\mathbf{Relint}}
\title{{{On determinantal variety}}}

\author{{{Li Enyu}}}

\date{}

\begin{document}

\maketitle

\def\H{\mathscr H}
\def\deter{\mathbb R_{\le r}^{m\times n}}
\section{The determinantal variety}The determinantal variety $\mathbb R_{\le r}^{m\times n}$ consists of matrices $\R^{m\times n}$ with rank less than or equal to $r$. The natural action of $O(m)\times O(n)$ gives an orbit decomposition \gs{
    \deter=\bigcup_{s\le r}\R_s^{m\times n}
} where each orbit is the smooth manifold consists of matrice of fixed rank (There is also an explicit cell decomposition but we do not need here).
Now  throuout this article let's fix a matrice with rank $s<r$ in an SVD coordinate \gs{
    X=\jz{
        \Lambda&0\\0&0
    }
}where $\Lambda$ is diagonal of order $r$. To visualize, note that a neighborhood of $X$ lies in $\bigcup_{\ge s}\R_s^{m\times n}$, namely there is a manifold part that preserving the rank, but also some direction that increasing rank, and the former forms the boundary of the latter. 
Consider the diffeomorphism near a neighborhood of $X$\gs{
\phi:U\to V\times \R_{\le r-s}^{(m-s)\times (n-s)}\subset\R^{s(m+n-s)}\times \R_{\le r-s}^{(m-s)\times (n-s)},\jz{A&B\\C&D}\mapsto (A,B,C| D/A)
}Here $D/A=D-CA^{-1}B$.
This implies locally $\R_{\le r}^{m\times n}$ \textbf{is a product of the manifold part and singular part}. Also note $d\phi$ is identity at $X$.

The tangent space $T_X\R^{m\times n}_s$ consists of matrices \gs{
    Y_1=\jz{*&*\\ *&0}
}(
Throughout this article $*$ means arbitrary matrices.)
But the total tangent cone in $\deter$ consists of \gs{
    Y_2=\jz{*&*\\ *&D}, \text{rank}{D}\le r-s
}

\section{The invariant manifold}
Now consider a constraint $h:\R^{m\times n}\to \R^k$ that has constant rank and is invariant under the right action of $O(n)$. This implies the regular level set $\H=h^{-1}(0)$ has following properties:
\begin{enumerate}
    \item $\H$ is invariant under action of $O(n)$. 

    \item The stablizer $G_X$ acts on $N_X\H$ by identity. 
\end{enumerate}
These are \textbf{the only qualifications} we need. (In fact for transverality, 2 is enough)
The action of $O(n)$ gives a fundamental vector field on $\R_s^{m\times n}$, at $X$ this gives following tangent vector:
\gs{
    Y_1=\jz{
        \Lambda&0\\0&0
    }\Omega_0= \jz{
        \Lambda \Omega&*\\0&0
    }, \oo_0\in \mathfrak{so}_n,\oo\in \mathfrak{so}_s
}(i.e. here $\Omega_0$ and $\oo$ are skew-symmetric matrices.)

Also note the stablizer \gs{
    \sigma=\jz{0&0\\0&-I}\in O(n)_X
}
will flip the vectors of the form \gs{
    Y_2=\jz{0&0\\0&*}
}but by our property (2), $\sigma$ acts on normal space by identity, so $Y_2$ is automatically in $T_XH$. 

Another way to see this is to consider the fibration\gs{
   \pi: \R^{m\times n}\to \text{sym}^+; X\mapsto XX^*
}
$\pi$ is not a submersion globally but it is when restricted to each fixed rank $s$ with fiber $O(s)$, so $h$ actually factors thru this map. Now it's not hard to see the vectors $Y_1+Y_2$ are exactly kernel of $d\pi$. (One can also use invariant theory to get this result, the factorization is actually globally smooth.)

\section{Transverse intersection}
Recall in section 1 we have a local diffeomorphism $\phi:\R^{m\times n} \subset V\times R^{(m-s)(n-s)}$ at $X$ that sends $U$ to $V\times \R^{(m-s)(n-s)}_{\le r}$ with $d_X\phi$ identity. Now we can flatten $\H$ in this coordinate. This is standard in geometry. For completeness we give a proof here.

The projection $f:\H\cap \yk{V\times R^{(m-s)(n-s)}}\to R^{(m-s)(n-s)}$ is a submersion at $X$. Let $p$ be a projection of $V$ onto $\ker (d_X f)=T_X\H\cap \yk{V\times 0}$ and $q$ be its complement. Consider \gs{
    \psi:\H\cap \yk{V\times R^{(m-s)(n-s)}}\to \ker d_Xf\times R^{(m-s)(n-s)}, (x,y)\mapsto (p(x),y)
}This is a chart of $\H$ s.t. its locally a product. Now we know that under this chart, 
\gs{
     \R_{\le r}^{mn}=& V\times \R^{(m-s)(n-s)}_{\le r}=\hk{\jz{*&*\\ *&D}:\text{rank}(D)\le r-s}\\ \H=&\H^s\times \R^{(m-s)(n-s)}=\hk{\jz{ A&* \\ B&*}:(A,B)\in \H^s}
}where $\H^s$ is some manifold in $\R^{s(m+n-s)}$. Now its clear  that their intersection is transverse as they are transeverse in first component and at second components is an inclusion $\R^{(m-s)(n-s)}_{\le r}\subset \R^{(m-s)(n-s)}$. As a consequence
\gs{
    T_X\yk{\R_{\le r}^{mn}\cap \H}=&T_X\yk{\R_{\le r}^{mn}}\cap T_X\H
}
This immediately gives the Theorem 1 in \cite{pap}
\begin{thebibliography}{9}

\bibitem{pap} Yang, Y., Gao, B. \& Yuan, Yx. A space-decoupling framework for optimization on bounded-rank matrices with orthogonally invariant constraints. Math. Program. (2026).

\end{thebibliography}

\end{document}